%-*-latex-*-
%-*-latex-*-
\newcommand\COURSE{ciss240}
\newcommand\ASSESSMENT{q1904}
\newcommand\ASSESSMENTTYPE{Quiz}
\newcommand\POINTS{\textwhite{xxx/xxx}}

\input{myquizpreamble}
\input{yliow}
\input{\COURSE}
%-*-latex-*-

%-*-latex-*-

%-*-latex-*-
\renewcommand\TITLE{\ASSESSMENTTYPE \ \ASSESSMENT}

\renewcommand\EMAIL{}
\renewcommand\AUTHOR{eagraham1@cougars.ccis.edu}

\textwidth=6in
\begin{document}
\topmatterthree

This is a closed-book, no compiler, 5 minute quiz.

%------------------------------------------------------------------------------
\nextq
The following program does not run. Insert one statement to correct the problem.
\\
\ANSWER
\begin{answercode}
#include <iostream>

// insert statement here
int sum(int start, int end, int step);

int main()
{
    std::cout << sum(5, 10, 1) << std::endl;
    return 0;
}

int sum(int start, int end, int step)
{
    int s = 0;
    for (int i = start; i <= end; i += step)
    {
        s += i;
    }
    return s;
}
\end{answercode}

%-------------------------------------------------------------------------------
\nextq
The following program does not run. Insert one statement to correct the problem.
\\
\ANSWER
\begin{answercode}
#include <iostream>
// insert statement here
void swap(int &x, int &y);

int main()
{
    int i, j;
    std::cin >> i >> j;
    swap(i, j);
    std::cout << i << ' ' << j << '\n';
    return 0;
}

void swap(int &x, int &y)
{
    int t = x;
    x = y;
    y = t;
}
\end{answercode}

%------------------------------------------------------------------------------
\nextq
The function \verb!play_audio()! plays an audio file. The parameters are:
\begin{enumerate}
\li \verb!filename!: a C-string that is the name of the audio file to be loaded.
\li \verb!track_number!: an integer.
$0$ is the first track, $1$ is the second track, etc.
\li \verb!repeat!: a boolean.
\end{enumerate}
The function returns $0$ if there are no errors,
returns $-1$ if the file cannot be found using the \verb!filename!,
and returns $-2$ if the \verb!track_number! is not found.
Write down the prototype of this function.
\\
\ANSWER
\begin{answercode}
int play_audio(const char* filename, int track_number, bool repeat);
\end{answercode}

%------------------------------------------------------------------------------
\nextq
The function \verb!copy_array()! copies the values of an array to another.
The parameters are:
\begin{enumerate}
\li \verb!x!: the target integer array.
\li \verb!x_length!: an integer reference variable.
\li \verb!y!: the source integer array
\li \verb!y_length!: an integer variable.
\end{enumerate}
When you call \texttt{copy\_array(x, x\_length, y, y\_length)} the values
\verb!y[0]!,...,\texttt{y[y\_length - 1]} are copied to
\verb!x[0]!,...,\verb!x[y_length - 1]! (respectively) and
\verb!x_length! is set to \verb!y_length!.
Write down the prototype of this function.
\\
\ANSWER
\begin{answercode}
void copy_array(int x[], int &x_length, const int y[], int y_length);
\end{answercode}

%------------------------------------------------------------------------------
\nextq
What is a suitable prototype for \verb!array_equals()! that
compares the values of two integer arrays \verb!x! and \verb!y!
from returns \verb!true!
when
\verb!x[0]!,...,\verb!x[n - 1]!
and
\verb!y[0]!,...,\verb!y[n - 1]!
are the same (respectively), and \verb!false! otherwise.
\\
\ANSWER
\begin{answercode}
bool array_equals(const int x[], const int y[], int n);
\end{answercode}

%------------------------------------------------------------------------------
%-*-latex-*-
\newpage
\input{instructions}
\end{document}

